Fast OT backends reduce the cost of one Sinkhorn update, but fixed-width
batched execution still runs at full width until its slowest problem finishes.
We present \policy{}, an execution layer for batched entropic optimal transport
that dynamically removes completed problems and narrows subsequent updates.
It starts from standard batched execution of independent Sinkhorn problems.
After an alternating scaling cycle, one marginal is current by construction;
a batched screening of the other marginal filters the problems that require the
two-marginal verifier. Problems that pass the stopping rule are removed from
every per-problem tensor, while unfinished
problems continue.

We develop a hardware-aware performance model. For first-passing iteration
numbers $n_1,\ldots,n_W$ and active batch width
$A_k=|\{t:n_t\geq k\}|$, static execution performs $Wn_{\max}$ problem
updates, whereas dynamic compaction performs exactly
$\sum_k A_k=\sum_t n_t$.
If $c_k(a,\xi_k)$ is the measured update cost at width $a$, the exact
matched-path runtime difference is the update-time reduction
$\sum_k[c_k(W,\xi_k)-c_k(A_k,\xi_k)]$ minus incremental control overhead.
Different completion depths create removable work, while measured kernel time
at each active width determines the runtime benefit. On the measured
A100/Triton configuration, wave quantization makes update time stepwise: for
$n=1024$, widths one through six have similar cost, whereas larger problems
respond to width reduction earlier.

Matched fixed-width/dynamic-compaction comparisons give
\CMetroCorrectedEndpointSpeedup{} on the four-worker MetroPT-3 endpoint,
$1.054\times$ for ImageNet-32 Sinkhorn time, and up to
$3.110\times$ in the MetroPT-3 scale-out study. An OTT-JAX realization
gives $1.366$--$1.581\times$ across three tolerances. When one initialization is
used for multiple problems, matched packed-Triton dynamic compaction gives $1.421\times$ and
removes 30.62\% of problem updates. In Packer19, a matched logical-masking
control isolates compaction, and the gain vanishes when all problems finish at
the same verification round. All reported consumer and training-quality checks
pass.
